\section{Conclusion}

This paper studies \emph{document-level visual retrieval}, which retrieves complete documents rather than isolated pages. To address this setting, we propose DocPC, a document-level visual retriever based on representative page composition, and introduce DocViRe, a benchmark with document-grounded queries and multi-positive annotations.

DocPC composes representative pages into compact $2{\times}2$ grids, enabling efficient indexing while preserving document-level semantic signals. Experiments on DocViRe show that DocPC achieves the best overall average performance among all compared methods. In particular, DocPC-ColQwen (First-4) reaches an average NDCG@5 of 44.09, outperforming the strongest page-level baseline ColQwen-page (DocViRe-sft, Boundary) at 38.91, while reducing indexed images, vectors, and storage footprint by 10.1$\times$.

Overall, the results show that document-level composition is more effective than page-level aggregation under the same supervision, and that document-level visual retrieval is both practical and promising for future research.

\section{Limitations}
This work has several limitations. First, while we explored both fixed and embedding-based adaptive page selection strategies, investigating more sophisticated adaptive allocation methods that dynamically adjust to document characteristics is a promising direction for future work. Second, DocViRe is constructed from English PDF documents across seven domains. Broader coverage of languages, document genres, and retrieval settings remains for future work. Third, the benchmark is designed for document-level retrieval and does not yet cover downstream settings such as interactive retrieval, multi-step search, or retrieval-augmented generation.